\section{Khớp mô hình thay thế thưa}
\label{sec:ch03-khop-mo-hinh-thay-the}

\index{hàm nhân}%
Các mẫu đã nhiễu không đóng góp ngang nhau. LIME dùng một \tn{hàm nhân}{kernel}
$\pi_x(z)$ để gán trọng số lớn hơn cho mẫu gần $x$. Với lớp mô hình tuyến tính,
bài báo dùng hàm mất mát bình phương có trọng số cục bộ và hàm nhân mũ theo một
hàm khoảng cách $D$ cùng độ rộng $\sigma$~\autocite{p09lime}:

\begin{equation}
\label{eq:ch03-mat-mat-cuc-bo}
\mathcal{L}(f,g,\pi_x)=
\sum_{(z,z',\,f(z),\,\pi_x(z)) \in \mathcal{Z}}
\pi_x(z)\bigl(f(z)-g(z')\bigr)^2,
\qquad
\pi_x(z)=\exp\left(-\frac{D(x,z)^2}{\sigma^2}\right).
\end{equation}

Ở đây $\mathcal{Z}$ là tập mẫu mà thuật toán dựng: mỗi phần tử mang đầu vào
$z$ đã được dựng lại, biểu diễn đã nhiễu $z'$ sinh ra nó, đầu ra $f(z)$ và
trọng số $\pi_x(z)$. Một mẫu xa vẫn có thể xuất hiện trong phép khớp, nhưng trọng
số của nó nhỏ hơn. Cách chọn $D$ phụ thuộc miền:
bài báo nêu khoảng cách cosine cho văn bản và khoảng cách Euclid cho ảnh. Như
vậy, từ \enquote{cục bộ} không chỉ là trực giác hình học. Nó được thực hiện bằng
một hàm khoảng cách, một hàm nhân và một độ rộng cụ thể.

\index{K-LASSO}%
Để mô hình còn đọc được, LIME giới hạn số hệ số khác không của $g$. Bài báo
chọn trước số lượng $K$ thành phần, dùng Lasso theo regularization path để
chọn chúng, rồi khớp lại các trọng số bằng bình phương tối thiểu. Quy trình đó
được bài báo gọi là K-LASSO~\autocite{p09lime}. Thuật toán~\ref{alg:ch03-lime}
viết lại quy trình này dưới dạng pseudo-code. Nó gộp thuật toán được đánh số
của bài báo với hai bước mà phần văn xuôi quanh đó mô tả, là việc dựng lại
đầu vào và việc trả về mô hình tuyến tính; nó không phải mã đã chạy trong
cuốn sách.

\begin{algorithm}[htbp]
  \caption{Pseudo-code cho lời giải thích LIME thưa, dựng theo thuật toán
  được đánh số của Ribeiro, Singh và Guestrin cùng phần văn xuôi mô tả quy
  trình quanh nó~\autocite{p09lime}}
  \label{alg:ch03-lime}
  \begin{algorithmic}[1]
    \Require mô hình hộp đen $f$, đầu vào $x$, biểu diễn $x'$, hàm nhân $\pi_x$,
    số mẫu $N$, độ dài lời giải thích $K$
    \State $\mathcal{Z} \gets \varnothing$
    \For{$i \gets 1$ \textbf{to} $N$}
      \State $z'_i \gets \textsc{SinhNhiễu}(x')$
      \State $z_i \gets \textsc{DựngLạiĐầuVào}(z'_i)$
      \State $\mathcal{Z} \gets \mathcal{Z} \cup
      \{(z_i,z'_i,f(z_i),\pi_x(z_i))\}$
    \EndFor
    \State $w \gets \textsc{K-LASSO}(\mathcal{Z},K)$
    \State \Return mô hình tuyến tính $g(z')=w\mathbin{\cdot}z'$
  \end{algorithmic}
\end{algorithm}

Quy trình ấy chạy hết trên một ví dụ đủ nhỏ để tính bằng tay, và cuốn sách quay
lại đúng ví dụ này ở ba chương sau. Lấy
\[
  f(x_1,x_2) = x_1 + 2x_2 + x_1 x_2 ,
\]
cần giải thích tại $x = (1,1)$. \tn{Biểu diễn có thể diễn giải}{interpretable
representation} là hai bit bật tắt
từng \tn{đặc trưng}{feature}, và dựng lại đầu vào nghĩa là đặt đặc trưng bị tắt về 0. Với hai
đặc trưng thì chỉ có bốn mẫu, nên ở đây tập $\mathcal{Z}$ là toàn bộ chúng chứ
không phải một mẫu ngẫu nhiên; đó là lý do ví dụ tính được hết.

Chọn $D$ là khoảng cách Euclid và $\sigma^2 = 1/\ln 2$, thì
$\pi_x(z) = \exp(-D^2/\sigma^2)$ nhận đúng các giá trị trong bảng dưới. Độ rộng
ấy được chọn cho số đẹp; nó là một cấu hình minh họa, không phải mặc định của
LIME.

\begin{table}[htbp]
  \centering
  \caption{Bốn mẫu của ví dụ, tính từ $f(x_1,x_2) = x_1 + 2x_2 + x_1x_2$ với
  $x = (1,1)$ và đặc trưng bị tắt đặt về 0. Bốn dòng là toàn bộ tập mẫu chứ
  không phải một mẫu rút ngẫu nhiên, và các số là phép tính tay chứ không phải
  kết quả chạy một thư viện nào.}
  \label{tab:ch03-vi-du-bon-mau}
  \begin{tabular}{@{}cccc@{}}
    \toprule
    $z'$ & $f(z)$ & $D(x,z)^2$ & $\pi_x(z)$ \\
    \midrule
    $(0,0)$ & $0$ & $2$ & $1/4$ \\
    $(1,0)$ & $1$ & $1$ & $1/2$ \\
    $(0,1)$ & $2$ & $1$ & $1/2$ \\
    $(1,1)$ & $4$ & $0$ & $1$ \\
    \bottomrule
  \end{tabular}
\end{table}

Lấy $K = 2$, tức giữ cả hai đặc trưng, nên bước chọn biến không loại gì và chỉ
còn bước khớp có trọng số. Mô hình thay thế là $g(z') = w_1 z'_1 + w_2 z'_2$,
không có hệ số tự do, đúng dạng mà thuật toán trả về. Mẫu $(0,0)$ không góp gì
vào phép khớp vì cả hai biến đều bằng không ở đó. Đạo hàm
$\mathcal{L}$ theo $w_1$ và theo $w_2$ rồi cho bằng không, sau khi nhân hai vế
với 2, cho hệ
\[
  3w_1 + 2w_2 = 9, \qquad 2w_1 + 3w_2 = 10 ,
\]
nên $w_1 = 7/5$ và $w_2 = 12/5$. Lời giải thích đọc được là: đặc trưng thứ hai
quan trọng hơn đặc trưng thứ nhất, theo tỉ lệ $12$ trên $7$.

Một chi tiết của kết quả ấy đáng dừng lại, vì nó là chỗ chương sau bắt đầu.
Tại chính điểm cần giải thích, $g(1,1) = 19/5 = 3.8$, trong khi $f(1,1) = 4$.
Mô hình thay thế không buộc phải trùng mô hình gốc ngay tại điểm được giải
thích, và ở đây nó không trùng: tổng bình phương sai số có trọng số còn lại
bằng $1/5$. Phép khớp tìm một mặt phẳng hợp nhất với bốn mẫu đã cân trọng số,
không tìm một mặt phẳng đi qua một điểm.

